% !TEX program = xelatex
% =====================================================================
% FUENTE DE VERDAD del MOF (decisión 2026-07-17)
% Editar este archivo y preamble/preambulo.tex (paquetes, colores de
% marca Yachay y metadatos del título).
% Compilar:  latexmk -xelatex -cd MOF/MOF_v1.tex
% (XeLaTeX es el motor oficial: fija las fuentes del render de Quarto)
% El antiguo MOF_v1.qmd fue eliminado (2026-07-17); se conserva en el
% historial de git si se necesita consultar.
% =====================================================================
% Options for packages loaded elsewhere
\PassOptionsToPackage{unicode}{hyperref}
\PassOptionsToPackage{hyphens}{url}
\PassOptionsToPackage{dvipsnames,svgnames,x11names}{xcolor}
%
\documentclass[
  spanish,
]{article}

% Preámbulo separado: paquetes, colores de marca Yachay y metadatos del título
% Editar en preamble/preambulo.tex
\usepackage{xcolor}
\usepackage[margin=1in]{geometry}
\usepackage{amsmath,amssymb}
\setcounter{secnumdepth}{5}
\usepackage{iftex}
\ifPDFTeX
  \usepackage[T1]{fontenc}
  \usepackage[utf8]{inputenc}
  \usepackage{textcomp} % provide euro and other symbols
\else % if luatex or xetex
  \usepackage{unicode-math} % this also loads fontspec
  \defaultfontfeatures{Scale=MatchLowercase}
  \defaultfontfeatures[\rmfamily]{Ligatures=TeX,Scale=1}
\fi
\usepackage{lmodern}
\ifPDFTeX\else
  % xetex/luatex font selection
\fi
% Use upquote if available, for straight quotes in verbatim environments
\IfFileExists{upquote.sty}{\usepackage{upquote}}{}
\IfFileExists{microtype.sty}{% use microtype if available
  \usepackage[]{microtype}
  \UseMicrotypeSet[protrusion]{basicmath} % disable protrusion for tt fonts
}{}
\makeatletter
\@ifundefined{KOMAClassName}{% if non-KOMA class
  \IfFileExists{parskip.sty}{%
    \usepackage{parskip}
  }{% else
    \setlength{\parindent}{0pt}
    \setlength{\parskip}{6pt plus 2pt minus 1pt}}
}{% if KOMA class
  \KOMAoptions{parskip=half}}
\makeatother
% Make \paragraph and \subparagraph free-standing
\makeatletter
\ifx\paragraph\undefined\else
  \let\oldparagraph\paragraph
  \renewcommand{\paragraph}{
    \@ifstar
      \xxxParagraphStar
      \xxxParagraphNoStar
  }
  \newcommand{\xxxParagraphStar}[1]{\oldparagraph*{#1}\mbox{}}
  \newcommand{\xxxParagraphNoStar}[1]{\oldparagraph{#1}\mbox{}}
\fi
\ifx\subparagraph\undefined\else
  \let\oldsubparagraph\subparagraph
  \renewcommand{\subparagraph}{
    \@ifstar
      \xxxSubParagraphStar
      \xxxSubParagraphNoStar
  }
  \newcommand{\xxxSubParagraphStar}[1]{\oldsubparagraph*{#1}\mbox{}}
  \newcommand{\xxxSubParagraphNoStar}[1]{\oldsubparagraph{#1}\mbox{}}
\fi
\makeatother


\usepackage{longtable,booktabs,array}
\usepackage{calc} % for calculating minipage widths
% Correct order of tables after \paragraph or \subparagraph
\usepackage{etoolbox}
\makeatletter
\patchcmd\longtable{\par}{\if@noskipsec\mbox{}\fi\par}{}{}
\makeatother
% Allow footnotes in longtable head/foot
\IfFileExists{footnotehyper.sty}{\usepackage{footnotehyper}}{\usepackage{footnote}}
\makesavenoteenv{longtable}
\usepackage{graphicx}
\makeatletter
\newsavebox\pandoc@box
\newcommand*\pandocbounded[1]{% scales image to fit in text height/width
  \sbox\pandoc@box{#1}%
  \Gscale@div\@tempa{\textheight}{\dimexpr\ht\pandoc@box+\dp\pandoc@box\relax}%
  \Gscale@div\@tempb{\linewidth}{\wd\pandoc@box}%
  \ifdim\@tempb\p@<\@tempa\p@\let\@tempa\@tempb\fi% select the smaller of both
  \ifdim\@tempa\p@<\p@\scalebox{\@tempa}{\usebox\pandoc@box}%
  \else\usebox{\pandoc@box}%
  \fi%
}
% Set default figure placement to htbp
\def\fps@figure{htbp}
\makeatother



\ifLuaTeX
\usepackage[bidi=basic]{babel}
\else
\usepackage[bidi=default]{babel}
\fi
% get rid of language-specific shorthands (see #6817):
\let\LanguageShortHands\languageshorthands
\def\languageshorthands#1{}


\setlength{\emergencystretch}{3em} % prevent overfull lines

\providecommand{\tightlist}{%
  \setlength{\itemsep}{0pt}\setlength{\parskip}{0pt}}



 


\usepackage{sectsty}
\definecolor{yachaypetroleo}{HTML}{1B4965}
\definecolor{yachayacero}{HTML}{5FA8D3}
\allsectionsfont{\color{yachaypetroleo}}
\makeatletter
\@ifpackageloaded{caption}{}{\usepackage{caption}}
\AtBeginDocument{%
\ifdefined\contentsname
  \renewcommand*\contentsname{Tabla de contenidos}
\else
  \newcommand\contentsname{Tabla de contenidos}
\fi
\ifdefined\listfigurename
  \renewcommand*\listfigurename{Listado de Figuras}
\else
  \newcommand\listfigurename{Listado de Figuras}
\fi
\ifdefined\listtablename
  \renewcommand*\listtablename{Listado de Tablas}
\else
  \newcommand\listtablename{Listado de Tablas}
\fi
\ifdefined\figurename
  \renewcommand*\figurename{Figura}
\else
  \newcommand\figurename{Figura}
\fi
\ifdefined\tablename
  \renewcommand*\tablename{Tabla}
\else
  \newcommand\tablename{Tabla}
\fi
}
\@ifpackageloaded{float}{}{\usepackage{float}}
\floatstyle{ruled}
\@ifundefined{c@chapter}{\newfloat{codelisting}{h}{lop}}{\newfloat{codelisting}{h}{lop}[chapter]}
\floatname{codelisting}{Listado}
\newcommand*\listoflistings{\listof{codelisting}{Listado de Listados}}
\makeatother
\makeatletter
\makeatother
\makeatletter
\@ifpackageloaded{caption}{}{\usepackage{caption}}
\@ifpackageloaded{subcaption}{}{\usepackage{subcaption}}
\makeatother
\usepackage{bookmark}
\IfFileExists{xurl.sty}{\usepackage{xurl}}{} % add URL line breaks if available
\urlstyle{same}
\hypersetup{
  pdftitle={Manual de Organización y Funciones},
  pdflang={es},
  colorlinks=true,
  linkcolor={yachaypetroleo},
  filecolor={Maroon},
  citecolor={Blue},
  urlcolor={yachayacero},
  pdfcreator={LaTeX via pandoc}}


\title{Manual de Organización y Funciones}
\usepackage{etoolbox}
\makeatletter
\providecommand{\subtitle}[1]{% add subtitle to \maketitle
  \apptocmd{\@title}{\par {\large #1 \par}}{}{}
}
\makeatother
\subtitle{Yachay Research Associates LLC}
\author{}
\date{2026-02-11}


\begin{document}
\maketitle

\renewcommand*\contentsname{Tabla de contenidos}
{
\hypersetup{linkcolor=yachaypetroleo}
\setcounter{tocdepth}{3}
\tableofcontents
}

\newpage

\begin{quote}
\textbf{Versión:} 2.0\\
\textbf{Elaborado por:} Darwin del Castillo\\
\textbf{Aprobado por:} Darwin del Castillo\\
\end{quote}

\newpage

\section{Introducción}\label{introducciuxf3n}

\subsection{Objetivo del manual}\label{objetivo-del-manual}

Establecer la estructura organizacional mínima, las responsabilidades y
funciones de cada cargo, así como los procesos de trabajo y políticas
internas para asegurar entregables reproducibles, calidad técnica y un
servicio consistente a clientes de salud (farma, biotech, CROs y,
progresivamente, ONGs) en Latinoamérica.

\subsection{Alcance}\label{alcance}

Aplica a: - Operaciones comerciales, administrativas y de entrega de
proyectos. - Proyectos de analítica de datos, \textbf{RWE},
diseño/ejecución de análisis, apoyo estadístico a \textbf{ensayos
clínicos}, y \textbf{economía de la salud (HEOR)}. - Colaboradores
internos y externos (contratistas) cuando existan.

\subsection{Marco legal (plantilla)}\label{marco-legal-plantilla}

\begin{quote}
Ajusta esta sección a tu realidad legal (EE. UU. / Perú) y asesoría
contable-jurídica.
\end{quote}

\begin{itemize}
\tightlist
\item
  Cumplimiento de contratos, confidencialidad y propiedad intelectual
  según el país del cliente y el país de contratación.
\item
  Normas y buenas prácticas de protección de datos aplicables según
  jurisdicción (por ejemplo, datos personales, datos sensibles de salud,
  transferencias internacionales y acuerdos de procesamiento).
\item
  Reglas internas para manejo de datos, seguridad y trazabilidad de
  entregables.
\end{itemize}

\section{Estructura organizacional}\label{estructura-organizacional}

\subsection{Organigrama (estructura mínima --- 3
personas)}\label{organigrama-estructura-muxednima-3-personas}

La estructura se organiza en \textbf{una dirección general} y
\textbf{dos funciones clave}: (a) desarrollo de negocio/marketing y (b)
ejecución técnica con aseguramiento de calidad.

\textbf{Organigrama (texto):}

\begin{itemize}
\tightlist
\item
  \textbf{Director General (CEO)}

  \begin{itemize}
  \tightlist
  \item
    \textbf{Director(a) de Marketing y Desarrollo de Negocio (CMO/BD)}
  \item
    \textbf{Consultor(a) Senior RWE/Epidemiología y Analítica (Delivery
    Lead)}
  \end{itemize}
\end{itemize}

\subsection{Niveles jerárquicos}\label{niveles-jeruxe1rquicos}

\begin{enumerate}
\def\labelenumi{\arabic{enumi}.}
\tightlist
\item
  \textbf{Dirección}: CEO\\
\item
  \textbf{Función comercial y relación con el cliente}: CMO/BD\\
\item
  \textbf{Función técnica y entrega}: Consultor Senior (Delivery Lead)
\end{enumerate}

\subsection{Principios de gobernanza operativa
(ligera)}\label{principios-de-gobernanza-operativa-ligera}

\begin{itemize}
\tightlist
\item
  Las decisiones de \textbf{alcance, precio, contratación y riesgo} se
  aprueban por el CEO.
\item
  El CMO/BD lidera \textbf{pipeline comercial}, propuestas y
  coordinación de reuniones con stakeholders.
\item
  El Consultor Senior lidera la \textbf{ejecución técnica}, QA/QC y
  consistencia metodológica.
\item
  Toda entrega formal al cliente pasa por \textbf{revisión} (CEO o
  Consultor Senior, según el tipo de entregable).
\end{itemize}

\section{Descripción de cargos y
funciones}\label{descripciuxf3n-de-cargos-y-funciones}

\begin{quote}
Para cada cargo: propósito, funciones, responsabilidades, autoridad,
indicadores y requisitos mínimos.
\end{quote}

\subsection{Director General (CEO)}\label{director-general-ceo}

\textbf{Objetivo del puesto}\\
Definir la estrategia, asegurar sostenibilidad financiera, supervisar la
calidad global y mantener relaciones clave con clientes y socios.

\textbf{Funciones principales} - Definir posicionamiento, oferta de
servicios y roadmap trimestral. - Aprobar propuestas, presupuestos,
contratos y \emph{Statements of Work} (SoW). - Gestionar riesgos
(legales, seguridad de datos, reputación, cronograma). - Supervisar
calidad: estándar de reproducibilidad, documentación y QA/QC. - Definir
y monitorear KPIs (ingresos, margen, recurrencia, tasa de conversión).

\textbf{Responsabilidades} - Asegurar el cumplimiento de políticas de
confidencialidad y seguridad. - Definir criterios de aceptación de
entregables y planes de mitigación. - Priorizar el portafolio de
proyectos y asignar recursos. - Representación institucional ante
clientes estratégicos.

\textbf{Autoridad / decisiones} - Firma final de contratos y NDAs (o
delegación formal). - Aprobación de cambios mayores de alcance
(\emph{change requests}). - Aprobación de contratación de terceros y
compras de software.

\textbf{Indicadores sugeridos (KPIs)} - \% de ingresos recurrentes
(retainers) - Margen bruto por proyecto - \emph{Sales cycle} promedio
(semanas) y tasa de conversión - Satisfacción del cliente (NPS o
encuesta breve post-proyecto)

\textbf{Requisitos mínimos} - Experiencia en investigación/analítica
aplicada en salud. - Capacidad de comunicación ejecutiva y gestión de
stakeholders. - Competencias en ética, confidencialidad y manejo de
datos sensibles.

\subsection{Director(a) de Marketing y Desarrollo de Negocio
(CMO/BD)}\label{directora-de-marketing-y-desarrollo-de-negocio-cmobd}

\textbf{Objetivo del puesto}\\
Generar demanda y oportunidades de venta, liderar la preparación de
propuestas y fortalecer la marca en el mercado regional.

\textbf{Funciones principales} - Prospección y gestión de pipeline:
lista de cuentas objetivo, \emph{outreach}, seguimiento. - Preparación
de material comercial: one-pagers, pitch deck, casos de estudio
anonimizados. - Redacción y coordinación de propuestas: alcance,
cronograma, supuestos, entregables. - Coordinación de reuniones con
clientes: agenda, actas, acuerdos y próximos pasos. - Gestión de
reputación: contenido técnico-divulgativo (sin información sensible).

\textbf{Responsabilidades} - Mantener el CRM (o registro simple)
actualizado: leads, estado, probabilidad, próximos pasos. - Asegurar
consistencia del mensaje y claridad del alcance en cada propuesta. -
Documentar acuerdos comerciales y traspasar al equipo técnico
(handover).

\textbf{Autoridad / decisiones} - Proponer precios y paquetes (sujeto a
aprobación del CEO). - Definir la estrategia de contenido y calendario
editorial.

\textbf{Indicadores sugeridos (KPIs)} - \# de leads calificados/mes - \#
de propuestas enviadas/mes y \% ganadas - Tiempo promedio de respuesta a
solicitudes - Crecimiento de red y \emph{engagement} (si aplica)

\textbf{Requisitos mínimos} - Experiencia en marketing B2B o desarrollo
de negocio en salud/tecnología. - Redacción profesional (propuestas,
SoW) y habilidades de negociación. - Comprensión básica de
HEOR/RWE/estadística (para vender con precisión).

\subsection{Consultor(a) Senior RWE/Epidemiología y Analítica (Delivery
Lead)}\label{consultora-senior-rweepidemiologuxeda-y-analuxedtica-delivery-lead}

\textbf{Objetivo del puesto}\\
Garantizar la entrega técnica con rigor, reproducibilidad y calidad,
traduciendo necesidades del cliente en análisis y productos finales
listos para toma de decisiones.

\textbf{Funciones principales} - Liderar el diseño analítico: preguntas,
estimandos, supuestos, plan analítico (mini-SAP). - Ejecución técnica:
R/Python/SQL, modelamiento HEOR y/o análisis RWE. - Preparar
entregables: reportes, tablas, figuras, anexos reproducibles y
presentaciones. - Implementar QA/QC: validaciones, revisión cruzada,
auditoría de supuestos. - Soporte a reuniones técnicas con cliente:
explicar metodología y resultados.

\textbf{Responsabilidades} - Mantener repositorios privados, control de
versiones y trazabilidad. - Documentar decisiones técnicas (por ejemplo,
supuestos de modelo, definiciones de cohortes). - Gestionar riesgos
técnicos (calidad de datos, sesgos, limitaciones) y comunicar temprano.

\textbf{Autoridad / decisiones} - Definir estándares de QA y criterios
de aceptación técnica. - Proponer mejoras de pipeline y automatización
(sujeto a priorización del CEO).

\textbf{Indicadores sugeridos (KPIs)} - \% entregables aceptados sin
retrabajo mayor - Tiempo de ciclo por sprint / hito - Hallazgos QA
detectados internamente vs reportados por cliente - Reutilización de
componentes (plantillas/código)

\textbf{Requisitos mínimos} - Formación sólida en
epidemiología/biostatística/HEOR o ciencia de datos en salud. - Dominio
de análisis reproducible (control de versiones, documentación). -
Experiencia con datos observacionales y/o modelos de decisión en salud.

\section{Mapa de procesos y flujos de
trabajo}\label{mapa-de-procesos-y-flujos-de-trabajo}

\begin{quote}
Los procesos se diseñan para una operación pequeña: claridad de roles,
entregables estándar y control de cambios.
\end{quote}

\subsection{Proceso 1 --- Captación y
preventa}\label{proceso-1-captaciuxf3n-y-preventa}

\textbf{Objetivo:} generar y convertir leads en propuestas aprobadas.\\
\textbf{Entradas:} cuentas objetivo, requerimientos preliminares,
RFPs.\\
\textbf{Salidas:} diagnóstico 1--2 páginas, propuesta + SoW, estimación
de costos y cronograma.

\textbf{Actividades} 1. Identificación de lead (CMO/BD) 2.
\emph{Discovery call} (CMO/BD + CEO/Consultor, según complejidad) 3.
Diagnóstico breve y opciones de paquete (CMO/BD) 4. Propuesta + SoW +
timeline (CMO/BD; revisión CEO/Consultor) 5. Negociación y cierre (CEO)

\subsection{Proceso 2 --- Inicio de proyecto
(kickoff)}\label{proceso-2-inicio-de-proyecto-kickoff}

\textbf{Objetivo:} alinear alcance, datos y criterios de aceptación.\\
\textbf{Entradas:} contrato/SoW, NDA, inventario de datos.\\
\textbf{Salidas:} plan analítico (mini-SAP), cronograma por hitos, plan
de comunicación.

\textbf{Artefactos} - Minuta de kickoff - Mini-SAP (estimandos,
supuestos, covariables, criterios) - Plan de entrega (hitos y
\emph{review gates})

\subsection{Proceso 3 --- Ejecución técnica
(delivery)}\label{proceso-3-ejecuciuxf3n-tuxe9cnica-delivery}

\textbf{Objetivo:} producir análisis y entregables según SoW.\\
\textbf{Entradas:} datos, mini-SAP, definiciones y supuestos.\\
\textbf{Salidas:} resultados reproducibles, reportes y/o modelos.

\textbf{Modalidad recomendada} - Sprints semanales con demo interna -
Revisión de supuestos antes de ``congelar'' resultados

\subsection{Proceso 4 --- Control de calidad
(QA/QC)}\label{proceso-4-control-de-calidad-qaqc}

\textbf{Objetivo:} reducir errores y asegurar consistencia
metodológica.\\
\textbf{Mínimo viable para 3 personas} - Revisión cruzada de outputs
(Consultor Senior \(\leftrightarrow\) CEO) cuando sea posible. - Checklist QA por tipo de
proyecto (RWE, HEOR, ensayos).

\textbf{Ejemplos de controles} - Validación de definiciones de cohortes
y conteos. - Revisión de supuestos y sensibilidad en HEOR. -
Reproducibilidad: correr pipeline desde cero en entorno limpio.

\subsection{Proceso 5 --- Entrega y
cierre}\label{proceso-5-entrega-y-cierre}

\textbf{Objetivo:} entregar productos finales y acordar siguientes
pasos.\\
\textbf{Salidas:} deck ejecutivo, reporte técnico, anexos,
código/documentación (según contrato).

\textbf{Cierre} - Reunión de resultados - Acta de aceptación - Lecciones
aprendidas (interna) - Propuesta de retainer o fase 2 (CMO/BD + CEO)

\section{Matriz RACI (3 personas)}\label{matriz-raci-3-personas}

\begin{longtable}[]{@{}lrrr@{}}
\toprule\noalign{}
Proceso / Actividad & CEO & CMO/BD & Consultor Senior \\
\midrule\noalign{}
\endhead
\bottomrule\noalign{}
\endlastfoot
Prospección y pipeline & A & R & C \\
Discovery call & A & R & C \\
Propuesta y SoW & A & R & C \\
Negociación / contrato & R/A & C & C \\
Kickoff y mini-SAP & A & C & R \\
Ejecución técnica & C & I & R/A \\
QA/QC & A & I & R \\
Entrega al cliente & A & C & R \\
Cierre y retroalimentación & A & R & C \\
\end{longtable}

\textbf{Leyenda:} R = Responsible, A = Accountable, C = Consulted, I =
Informed

\section{Políticas y normas
internas}\label{poluxedticas-y-normas-internas}

\subsection{Confidencialidad y manejo de
datos}\label{confidencialidad-y-manejo-de-datos}

\begin{itemize}
\tightlist
\item
  Todo proyecto requiere NDA/claúsulas de confidencialidad.
\item
  Minimización de datos: usar solo lo necesario y con acceso
  restringido.
\item
  Repositorios privados, control de acceso y cifrado de información
  sensible.
\item
  Prohibición de usar datos de cliente para demos sin autorización
  explícita.
\end{itemize}

\subsection{Gestión de documentación y
reproducibilidad}\label{gestiuxf3n-de-documentaciuxf3n-y-reproducibilidad}

\begin{itemize}
\tightlist
\item
  Todo proyecto debe incluir:

  \begin{itemize}
  \tightlist
  \item
    diccionario de datos (cuando aplique),
  \item
    README de ejecución,
  \item
    control de versiones,
  \item
    registro de supuestos y decisiones.
  \end{itemize}
\item
  Entregables con fechas, versión y responsable.
\end{itemize}

\subsection{Control de cambios (scope
creep)}\label{control-de-cambios-scope-creep}

\begin{itemize}
\tightlist
\item
  Cualquier cambio de alcance se gestiona mediante:

  \begin{itemize}
  \tightlist
  \item
    solicitud de cambio,
  \item
    evaluación de impacto (tiempo/costo),
  \item
    aprobación del CEO,
  \item
    actualización del SoW.
  \end{itemize}
\end{itemize}

\subsection{Capacitación y
desarrollo}\label{capacitaciuxf3n-y-desarrollo}

\begin{itemize}
\tightlist
\item
  Revisión mensual interna de buenas prácticas (QA, reproducibilidad,
  seguridad).
\item
  Librería interna de plantillas (reportes, tablas, figuras, decks).
\item
  Lecciones aprendidas por proyecto.
\end{itemize}

\section{Anexos (plantillas
sugeridas)}\label{anexos-plantillas-sugeridas}

\subsection{Anexo A --- Formato de informe (estructura
mínima)}\label{anexo-a-formato-de-informe-estructura-muxednima}

\begin{enumerate}
\def\labelenumi{\arabic{enumi}.}
\tightlist
\item
  Resumen ejecutivo\\
\item
  Contexto y objetivo\\
\item
  Datos y definiciones\\
\item
  Métodos\\
\item
  Resultados\\
\item
  Sensibilidades / limitaciones\\
\item
  Conclusiones y recomendaciones\\
\item
  Anexos reproducibles (tablas/figuras/código, según contrato)
\end{enumerate}

\subsection{Anexo B --- Matriz de aceptación de cambios (change
request)}\label{anexo-b-matriz-de-aceptaciuxf3n-de-cambios-change-request}

\begin{itemize}
\tightlist
\item
  Cambio solicitado
\item
  Motivo
\item
  Impacto en tiempo/costo
\item
  Aprobación (fecha, responsable)
\item
  Actualización del SoW
\end{itemize}

\subsection{Anexo C --- Lista de normativas aplicables
(placeholder)}\label{anexo-c-lista-de-normativas-aplicables-placeholder}

\begin{itemize}
\tightlist
\item
  Protección de datos (jurisdicción cliente)
\item
  Seguridad de información interna
\item
  Ética y cumplimiento en investigación/consultoría en salud
\end{itemize}

\section{Historial de revisión}\label{historial-de-revisiuxf3n}

\begin{longtable}[]{@{}llll@{}}
\toprule\noalign{}
Versión & Fecha & Descripción & Autor \\
\midrule\noalign{}
\endhead
\bottomrule\noalign{}
\endlastfoot
1.0 & 2026-02-11 & Versión inicial & Darwin del Castillo \\
\end{longtable}




\end{document}
